\section{Intersective and van der Corput Properties } \label{sec;intersective}
We turn to more explicitly arithmetic properties of these Cantor sets. Namely, the properties of being \emph{intersective} or \emph{van der Corput} sets.  

\begin{definition}
    We say that a set of integers $H \subset \mathbb N $ is \emph{intersective} if for 
 \begin{equation} \label{eSmax}
     I (H,N) = \max 
     \{ \lvert A\rvert /N \colon A\subset  [N], \ (A-A) \cap (H\setminus\{0\}) = \emptyset\}, 
 \end{equation}
 we have $I (H,N) \to 0$ as $N\to \infty$.   
 (The element $0$ must be excluded, since $0\in A-A$ for every non-empty $A$, and our Cantor sets contain $0$.) 
 We say that $H$ has a \emph{power savings (of order $c$)} if $I (H,N)  \ll N ^{-c}$.
 \end{definition}

It is the Theorem of S\'ark\"ozy \cite{MR487031} and Furstenberg \cite{MR0498471} that the square integers are intersective. 
S\'ark\"ozy also showed that the primes, plus or minus $1$, are also intersective. 
Both results have been the subject of a great deal of interest.  The best current result on the square integers is due to Green and Sawhney \cite{green2025newboundsfurstenbergsarkozytheorem}.   
 

We recall the definition of \emph{van der Corput set} of Kamae and Mend\`es-France
\cite{MR0516154}.  We give a quantitative definition  of this condition here. 

\begin{definition} 
For a set of integers $H \subset \mathbb{N}$  and $N\in \mathbb{N}$, set  
\begin{equation}\label{e;THN}
    \tau(H,N) = \inf _{\Psi \in \mathcal{T}_{N,H}} a_0
\end{equation}
where $\mathcal{T}_{N,H}$ is the set of all trigonometric polynomials $\Psi(x)$ satisfying 
\begin{equation} \label{e:T}
    \Psi(x) = a_0 + \sum_{ n<N, \ n\in H\setminus\{0\}} a_n \cos(2\pi n x) , 
\end{equation}
with $a_n\in \mathbb R$, $\Psi(0)=1$ and $\Psi(x)\geq 0$ for all $x$. 
We reserve $T$ for a measure preserving transformation, and use $\tau$ for the quantity \eqref{e;THN}. 
The sequence $H$ is said to be \emph{van der Corput (vdC)} if $\tau(H,N) \to 0$. 
That this agrees with the original definition of Kamae and Mend\`es-France, recalled below, is due to Ruzsa \cite{Ruzsa1984}. 
It is said to have a \emph{power savings of order $c$} in the van der Corput property if $\tau(H,N) \ll N ^{-c}$.  
\end{definition}


We recall some facts about the van der Corput property from \cite{MR0516154}: 
\begin{itemize}
    \item  $H$ is vdC if and only if it has the following property: for every sequence $\{x_n\} \subset [0,1)$, 
    if the sequences $\{x_{n+h}-x_n \colon n\in \mathbb N\}$ are uniformly distributed (mod $1$) for all $h\in H$, then $\{x_n\}$ is uniformly distributed (mod $1$). 
    This is the original definition of Kamae and Mend\`es-France. 

    \item A sufficient condition for $H$ to be vdC is that for every $n$ there are integers  $a_1 < a_2 < \cdots < a_n$ 
    with $a_j - a_i \in H$ for all $1\leq i < j \leq n$.  
    
    \item Kamae and Mend\`es-France \cite{MR0516154} showed that the van der Corput property for $H$ implies that $H$ is intersective. 
    In the quantitative form used here, $I(H,N) \ll \tau(H,N)$; see Ruzsa \cite{Ruzsa1984}. 
    The reverse implication does not hold, as was shown by Bourgain \cite{Bourgain1987}.  
\end{itemize}

We recall the definition of an IP-set, as defined by Furstenberg and Katznelson \cite{MR833409}: these are sets of the form $\{n_{t_1}+ \cdots +n_{t_k}\colon k\geq 1,\ t_1<t_2<\cdots<t_k\}$, that is, all possible finite sums of distinct elements of some infinite sequence $\{n_t\}$ of positive integers.
For a general digit set $D$, the Cantor set $\mathcal C_{b,D}$ contains an IP-set if and only if $0$ is an allowed digit; this is one reason for our standing assumption \eqref{e:standing}. 
Indeed, if $0\in D$, then for any $d\in D\setminus \{0\}$ the collection of sums $\sum_{t\in F}db^t$, over finite non-empty $F\subset \mathbb{N}_0$, is an IP-set contained in $\mathcal C$. 
Conversely, an IP-set generated by $\{n_t\}$ contains a non-zero multiple of every integer $q\geq 1$: among the partial sums $S_0=0$ and $S_v=n_1+\cdots +n_v$, two of $S_0,\ldots ,S_q$ agree modulo $q$, say $S_v\equiv S_{v'}$ with $v<v'$, and then $n_{v+1}+\cdots +n_{v'}$ lies in the IP-set and is divisible by $q$. 
Taking $q=b$, a non-zero element of $\mathcal C$ divisible by $b$ must have $0$ as its last digit. 
Under the standing assumption, $\mathcal C$ is vdC, since $\mathcal C$ contains an infinite sequence $u_0 < u_1 < \cdots$ with $u_j-u_i\in \mathcal C$, for $i<j$. 
Indeed, fix $d\in D\setminus\{0\}$ and set 
\begin{equation} \label{e:uj}
    u_j=\sum_{t<j}db^t , \qquad j\geq 0. 
\end{equation}
Thus, $\mathcal C$ is intersective. 
Making this argument quantitative is immediate, and gives a bound of order $1/\log N$.

\begin{proposition}\label{p:logSavings}
Let $\mathcal C=\mathcal C_{b,D}$ satisfy the standing assumption \eqref{e:standing}. Then, for all integers $N\geq 2$, 
\begin{equation}\label{e:logSavings}
\tau(\mathcal C, N) \leq \frac{\log b}{\log N} , \qquad\textup{and consequently}\qquad I(\mathcal C,N) \ll \frac{\log b}{\log N}. 
\end{equation}
\end{proposition}

\begin{proof}
Let $n\geq 1$ be an integer, and let $u_0<u_1<\cdots <u_{n-1}$ be as in \eqref{e:uj}, so that $u_0=0$. 
For $i<j$ we have $u_j-u_i=\sum_{i\leq t<j}db^t\in\mathcal C\setminus\{0\}$, and distinct pairs $i<j$ give distinct digit strings, hence distinct differences. 
Also, $u_{n-1}=d\frac{b^{n-1}-1}{b-1}\leq b^{n-1}-1$, so all these differences are less than $N_0\coloneqq b^{n-1}$. 
Consider the Fej\'er type kernel 
\begin{align}
    P (x) = \Bigl\lvert \Expectation_{0\leq j<n} e(u_jx)\Bigr\rvert^2 
    = \frac {1}{n ^2 } \sum_{i,j=0} ^{n-1} e\bigl( (u_j-u_i)x\bigr) . 
\end{align}
By construction $P\geq 0$ and $P(0)=1$. The frequencies $u_j-u_i$ and $u_i-u_j$ occur with the same coefficient $n^{-2}$, so $P$ is a cosine polynomial with frequencies in $\mathcal C\cap [0,N_0)$. 
Only the $n$ diagonal terms contribute to the constant term, so $a_0=n^{-1}$ and $\tau(\mathcal C,N_0)\leq n ^{-1} $. 
Given $N\geq 2$, take $n=1+\lfloor \log N/\log b\rfloor$, so that $N_0=b^{n-1}\leq N$ and $n\geq \log N/\log b$. 
As $N\mapsto \tau(\mathcal C,N)$ is nonincreasing, $\tau(\mathcal C,N)\leq \tau(\mathcal C,N_0)\leq n^{-1}\leq \log b/\log N$. 
The bound on $I(\mathcal C,N)$ follows from $I(\mathcal C,N)\ll \tau(\mathcal C,N)$. 
\end{proof}

We do not know whether every such Cantor set has power savings. 

\begin{question}\label{q:powersavings}
Let $\mathcal C=\mathcal C_{b,D}$ satisfy the standing assumption \eqref{e:standing}. 
Can \eqref{e:logSavings} be improved to a power savings, that is, is $\tau(\mathcal C,N)\ll N^{-c}$, or $I(\mathcal C,N)\ll N^{-c}$, for some $c=c(b,D)>0$? 
This is open to us even for the middle third Cantor set $\mathcal C_{3,\{0,2\}}$. 
Short of that, it would be of interest to improve the dependence on $b$ in \eqref{e:logSavings}. 
\end{question}
Power savings for Cantor sets with signed digits are established in \S\ref{sub:signed}. 
A necessary condition for a set of integers $H$ to be intersective is that $H$ contains a non-zero multiple of $q$, for every integer $q\geq 1$. 
(Otherwise, the sets $A=q\mathbb{Z}\cap[N]$ have $(A-A)\cap (H\setminus\{0\})=\emptyset$ and $\lvert A\rvert\geq N/q-1$.) 
In particular, if $0\notin D$, then $\mathcal C_{b,D}$ contains no multiple of $b$, and so is not intersective. 
 

 
\subsection{Power savings for signed digit sets} \label{sub:signed}
Green \cite{green2023sarkozystheoremshiftedprimes} has shown that the shifted primes have a power savings in the van der Corput property.  
We see below that a natural condition on the Cantor set implies that the  integer Cantor sets 
have a power savings. 

\emph{In this subsection only, we depart from the standing assumption \eqref{e:standing}.} 
Let $R_b\subset (-b,b)$ be a complete system of residues modulo $b$ with $0\in R_b$, and let $D\subset R_b$ with $0\in D$ and $2\leq \lvert D\rvert<b$. 
The digits in $D$ are allowed to be negative, but we insist that $\lvert d\rvert\leq b-1$ for $d\in D$. 
This normalization is used below to bound the elements of $\mathcal C_{b,D}$ in terms of the number of their digits: 
if $k\in\mathcal C_{b,D}$ has digits in the places $0,\ldots ,J$, then $\lvert k\rvert \leq (b-1)\sum_{j=0}^{J} b ^{j} <b ^{J+1} $. 
We define $\mathcal C_{b,D}$ by the same formula as before; the representation of an element of $\mathcal C_{b,D}$ by its digits is again unique. 
If $D$ contains negative digits, $\mathcal C_{b,D}$ contains negative integers. 
Since the intersective and van der Corput properties only concern the set of differences, which is symmetric, 
for such sets we apply the definitions above to $H=\{\lvert k\rvert \colon k\in \mathcal C_{b,D}\}\setminus\{0\}$. 
For instance, one can take $b=5$, $R_5=\{-2,-1,0,1,2\}$ and $D=\{-1,0,1\}$. 

\begin{proposition}\label{Powersavings}  
Let $\mathcal{C} = \mathcal{C}_{b,D}$ be a Cantor set with base $b$, and signed digit set $D\subset R_b$ as above.  
If $S$ is a finite set of integers with $S-S \subseteq D$ and $\lvert S\rvert >1$, then 
$\mathcal{C}$ is  van der Corput (hence intersective) with  power savings with $c =  \frac{ \log \lvert S\rvert }{\log b} $. 
    \end{proposition}

These examples supplement the primes minus one as concrete examples of sets of integers with a power savings in the intersective property, while also being `sparse', in that their density also decreases with a power savings.  

\begin{proof}
 Let $\sigma=\lvert S\rvert$, and note that $0\in S-S$ is an allowed digit. 
 For an integer $J\geq 0$, consider the product of Fej\'er type kernels 
 \begin{align}
     P_J (x) = \prod _{j=0}^J  \Bigl\lvert \Expectation_{s\in S} e(sb^jx)\Bigr\rvert^2 . 
 \end{align}
By construction, $P_J(0)=1$ and $P_J\geq 0$. 
Expanding each factor as $\sigma^{-2}  \sum_{s,t\in S} e((s-t)b^jx)$, we have 
\begin{align}
    P_J (x) = \sigma^{-2(J+1)} \sum_{ (s_j),(t_j)\in S^{J+1} } e\Bigl( 
    x\sum_{j=0}^J (s_j-t_j) b^j  \Bigr). 
\end{align}
Each frequency $n=\sum_{j=0}^J (s_j-t_j) b^j$ has digits $s_j-t_j\in S-S\subseteq D$, so $n\in\mathcal C$, and $\lvert n\rvert<b^{J+1}$, since $s_j-t_j\in D\subset R_b\subset (-b,b)$. 
As $S-S$ is symmetric, the frequencies $n$ and $-n$ appear with the same coefficient, so $P_J$ is a cosine polynomial, with real coefficients and frequencies in 
$\{\lvert k\rvert\colon k\in\mathcal C\}\cap[0,N)$, where $N=b^{J+1}$. 

It remains to compute the constant term, which is 
\begin{align}
    a_0 = \sigma^{-2(J+1)} \Bigl\lvert \Bigl\{  (s_j),(t_j)\in S^{J+1} \colon  \sum_{j=0}^J (s_j-t_j) b^j =0 \Bigr\}\Bigr\rvert. 
\end{align}
Since $s_j-t_j\in D\subset R_b$, and the representation of an element of $\mathcal C$ by its digits is unique, 
$\sum_{j=0}^J (s_j-t_j) b^j =0$ holds only if $s_j=t_j$ for all $j$. 
So only the $\sigma^{J+1}$ diagonal terms contribute, and 
\begin{align}
    a_0 = \sigma^{-2(J+1) + (J+1)} =\sigma^{-(J+1)}= N^{-\log\sigma/\log b}.   
\end{align}
This shows that $\tau(\mathcal C,N) \leq  N^{-\log\sigma/\log b}$ along the sequence  $N =b^{J+1}$, $J\geq0$. 
The function $N\mapsto \tau(\mathcal C,N)$ is nonincreasing, hence $\tau(\mathcal C,N) \leq b^{\log\sigma/\log b}  N^{-\log\sigma/\log b}=\sigma N^{-\log\sigma/\log b}$ 
for all integers $N\geq b$. 
\end{proof}

The above proposition is an instance of the following general principle. 
Let $E$ be a finite set of integers, and suppose that every non-zero difference of elements of $E$ lies in $\pm K$, for a set of positive integers $K$. 
If $A\subset[N]$ satisfies $(A-A)\cap K=\emptyset$, then the translates $A+e$, $e\in E$, are pairwise disjoint, and contained in an interval of length $N+\max E-\min E$. 
Hence $\lvert A\rvert\leq (N+\max E-\min E)/\lvert E\rvert$. 
In the Proposition, one takes $E=\{\sum_{j=0}^J s_jb^j\colon s_j\in S\}$; these sums are distinct, by uniqueness of the digit representation, so that 
$\lvert E\rvert=\lvert S\rvert^{J+1}$, and $\max E-\min E<b^{J+1}$ since the differences have digits in $D\subset (-b,b)$. 
Proposition \ref{p:logSavings} is the instance of the same principle in which $E=\{u_0,\ldots ,u_{n-1}\}$, with $u_j$ as in \eqref{e:uj}. 
In general, difference sets have a strong vdC property; see Example 2 in \cite{MR0516154}.

Here, we note simple sufficient conditions for $D$ to contain a difference set. 


\begin{proposition} \label{p:signedSufficient}
  Let $\mathcal C=\mathcal C_{b,D}$ be a Cantor set with signed digit set $D\subset R_b$ as above, and suppose either 
  \begin{enumerate}
      \item $D$ contains $d$ and $-d$ as integers, for some $d\neq 0$. 
      \item The residue system $R_b$ is symmetric ($R_b= \{-\frac {b-1}2,\ldots ,\frac {b-1}2\}$ for $b$ odd, 
      and $R_b=\{-b/2+1,\ldots b/2\}$ for $b$ even) and $\lvert D\rvert \geq  \lfloor b/2\rfloor+2$. 
  \end{enumerate}
Then, $\mathcal C$ is van der Corput with power savings at least $\log 2/\log b$. 

%     Suppose that $|D|> \frac{b+1}{2}$, or that $D$ contains a pair of antipodal points $d,-d\in \mathbb{Z}_b$. Then the Cantor set $C(b,D) $ is van der Corput with  power-savings of at least $\frac{\log 2}{\log b}$. 
\end{proposition} 

\begin{proof}
 If the first condition holds, the van der Corput property with power savings follows from Proposition \ref{Powersavings}. 
 And, the second condition implies the first: For $b$ odd, the residue system $R_b$ is partitioned into $0$ and $\frac {b-1}2$ 
 symmetric residue systems of the form $\{d,-d\}$ for $d=1,\ldots ,\frac {b-1}2$. That is a total of $\frac{b+1}2= \lfloor b/2\rfloor +1 $ residue systems.  
 And if $b$ is even, the residues $R_b \setminus \{0,b/2\}$, are partitioned into $\frac{b-2}2 = b/2-1$ symmetric residue systems. 
 In both cases the cardinality condition $\lvert D\rvert \geq  \lfloor b/2\rfloor+2$ implies that $D$ must contain a symmetric pair  of non-zero residues. 
   
    
%    Observe that if $|D|> \frac{b+1}{2}$, then $D$ contains a pair of nonzero antipodal points, so it suffices to prove that the claim follows from the second point. To do this, simply choose $S=\{0,d\}\subset \mathbb{Z}_b$ and apply Proposition \ref{Powersavings}.
\end{proof}

Next, we can give examples that relate the dimension of the Cantor set to the power savings.  In particular, there are small dimension Cantor sets that have power savings. 
(The shifted primes are full dimensional.) 


\begin{proposition} \label{p:dimension}
    For each  $\theta \in (0,1)$ and $ 0<\epsilon<\min\{\theta ,1-\theta \}$, there exists an integer Cantor set $\mathcal{C}=\mathcal{C}_{b,D}$, with a signed digit set $D$ as above, for which the following two points hold simultaneously: 
    \begin{enumerate}
        \item[(1)] $  \theta -\epsilon \leq \dim(\mathcal{C})=\frac{\log |D|}{\log b}  \leq \theta + \epsilon$, where $\dim(\mathcal C)=\lim_{N\to\infty}\frac{\log\lvert\mathcal C\cap[-N,N]\rvert}{\log N}$. 
        \item[(2)] There is a power savings of at least $\theta -\epsilon$ 
        for $\mathcal{C}$. 
    \end{enumerate}
\end{proposition}

\begin{proof}
First choose an odd base $b$ so large that 
\begin{equation} \label{e:bLarge}
    \frac{\log 5}{\log b}<\frac \epsilon 2 \qquad\textup{and}\qquad b ^{1-\epsilon /2}\leq b-1 , 
\end{equation}
and equip it with the symmetric residue system $R_b=\{-\frac{b-1}2,\ldots,\frac{b-1}2\}$. 
As $k$ ranges over $0,1,2,\ldots$, the quantity $\frac{\log (4k+1)}{\log b}$ increases from $0$ to $\infty$ in steps of at most $\frac{\log 5}{\log b}<\frac \epsilon2$. 
So we may fix $k\geq 1$ with $\frac{\log(4k+1)}{\log b}\in (\theta -\frac \epsilon 2,\ \theta +\frac \epsilon 2)$. 
Since $\epsilon <1-\theta $, we have $\theta +\frac \epsilon 2\leq 1-\frac \epsilon 2$, and so $4k+1<b ^{1-\epsilon /2}\leq b-1$ by \eqref{e:bLarge}. 
Therefore $D:=\{-2k,-2k+1,\ldots ,0,\ldots ,2k-1,2k\}$ is contained in $R_b$, and $\dim (\mathcal{C})=\frac{\log (4k+1)}{\log b}$, which is the point (1). 

For (2), put $S:=\{-k,\ldots ,0,\ldots ,k\}$, so that $S-S\subset D$. 
Proposition \ref{Powersavings} applies, and gives a power savings of 
\begin{equation}
\frac{\log (2k+1)}{\log b}=\dim \mathcal{C}-\frac{1}{\log b}\log \frac{4k+1}{2k+1} 
> \dim \mathcal{C}-\frac{\log 2}{\log b} > \Bigl(\theta -\frac \epsilon 2\Bigr)-\frac \epsilon 2 = \theta -\epsilon . 
\end{equation}
\end{proof}

\begin{remark}
The power savings of a set $H$ cannot exceed its dimension. 
Indeed, for any set $H$ of integers, choose $A\subset[N]$ greedily: run through $x=1,\ldots,N$, and add $x$ to $A$ if $x\notin a\pm (H\setminus\{0\})$ for every $a$ already in $A$. 
Each element of $A$ excludes at most $2\lvert H\cap[N]\rvert$ further elements, so $(A-A)\cap (H\setminus\{0\})=\emptyset$ and 
$I(H,N)\geq (2\lvert H\cap[N]\rvert+1)^{-1}$. 
For $H=\{\lvert k\rvert\colon k\in\mathcal C\}$ we have $\lvert H\cap[N]\rvert=N^{\dim(\mathcal C)+o(1)}$, so a power savings of order $c$ forces $c\leq\dim(\mathcal C)$. 
Since $I(H,N)\ll \tau(H,N)$, the same holds for the van der Corput property. 
Thus Proposition \ref{p:dimension} is sharp, up to $\epsilon$. 
\end{remark}




\subsection{Polynomial recurrence along Cantor sets}
From now on, we return to the standing assumption \eqref{e:standing}. 
We show that the intersective property continues to hold for certain polynomials evaluated along a Cantor set.  
The proof of Theorem \ref{t:polysIntersective} below draws on \S\ref{sec:uniformDistributionMod1} and \S\ref{sec:modular}, 
specifically on Theorem \ref{t:polyUniform}, Lemma \ref{l:Allq} and Lemma \ref{lem:positiveDensity}. 
Those sections do not depend on this one, so there is no circularity.  



\begin{theorem} \label{t:polysIntersective}
 Let $\mathcal C = \mathcal C_{b,D} = \{0=k_0<k_1< \cdots \}$, and $p $ be  a  polynomial 
 mapping the integers to the integers, with  positive leading coefficient 
 \emph{and no constant term.}
Set 
\begin{equation} \label{e:Hp}
H_p \coloneqq \{ \lvert p(k)\rvert \colon k \in \mathcal C\}\setminus \{0\} \subset \mathbb{N}. 
\end{equation}
Then, $H_p $ is intersective.  
Moreover, for any set of integers $A$ of positive upper Banach density $d^\ast(A)>0$, 
\begin{equation} \label{e;CantorDensity}
    \liminf_{N\to\infty} \Expectation_{0\leq n<N}
     d^\ast (A\cap (A - p(k_n)))  >0. 
\end{equation}
Consequently, there is a $c = c(A,p,\mathcal C)>0$ so that  the set of integers $\{n \colon d^\ast(A\cap (A-p(k_n))) >c \}$ 
has lower density at least $c$. 
\end{theorem}

The passage to absolute values in \eqref{e:Hp} is a convenience. 
Our definition of intersective is stated for subsets of $\mathbb{N}$, while a polynomial with positive leading coefficient and no constant term need be neither non-negative nor injective on $\mathcal C$: 
$p(x)=x^2-5x$ is negative at $x=1$, and $p(x)=x^2-x$ vanishes at $x=1$. 
Since the intersective property only concerns the symmetric set $A-A$, replacing $p(\mathcal C)$ by $H_p$ loses nothing.

A Cantor set contains an IP set, hence the recurrence property, namely that $d^\ast(A\cap A-p(k))>0$ for some non-zero $k\in\mathcal C$, is an immediate consequence of the polynomial IP recurrence theorem of Bergelson, Furstenberg and McCutcheon \cite{MR1417769}. 
Their result, and its multiple polynomial recurrence extension by Bergelson and McCutcheon \cite{BergelsonMcCutcheon2000}, show that the set of return times is large in the sense of IP sets: it is an IP$^\ast$ set. 
As far as we are aware, these results do not give information about the density of the return times measured along the enumeration $k_0<k_1<\cdots$ of $\mathcal C$. 
An IP$^\ast$ set must meet every IP set contained in $\mathcal C$, but this does not by itself imply that it contains a positive proportion of the $k_n$. 

The new information is the density result \eqref{e;CantorDensity}: The set of $n$ for which one has recurrence along $p(k_n)$ has positive lower density.  

For the proof, we rely upon the Furstenberg correspondence, the structure of dynamical systems, and  the following property of Cantor sets detailed in 
Lemma \ref{l:Allq}.  
In particular, that Lemma proves that for Cantor sets 
 $\mathcal C = \mathcal C_{b,D} 
= \{k_0<k_1 < \cdots\}$, 
the sequence $ k_n\mod q$ converges to a distribution on $\mathbb{Z} _q$, for all $q$, which assigns positive mass to $0$ (Lemma \ref{lem:positiveDensity}).  But   the limiting measures   are not always uniform. 


 


    
\begin{proof}
[Proof of Theorem \ref{t:polysIntersective}]
We use the Furstenberg correspondence principle \cite{MR0498471}. 
This gives us an \emph{invertible} system $(X,\mathcal B, \mu, T)$ --- it is the shift on a closed shift-invariant subset of $\{0,1\}^{\mathbb{Z}}$ --- 
and $\tilde A\subset X$ with $\mu(\tilde A)=d^\ast (A)$ for which 
$d^\ast(A\cap (A-m)) \geq \mu (\tilde A\cap T^{-m} \tilde A)$ for every integer $m$. 
And so, it suffices to show that 
\begin{equation} \label{eENA}
\liminf_N \Expectation _{0\leq n<N} \int_X \mathbf{1}_{\tilde A} T^{p(k_n)}  \mathbf{1}_{\tilde A} \;d\mu > 0. 
\end{equation}
In so doing, we will appeal to the fact that all terms are non-negative. In particular, we can always pass to averages over sets of $n$ of positive density. 


Split the Hilbert space $L^2(X)$ into $ \mathcal H _{\textup{rat}} \oplus \mathcal H _{\textup{0}}$, where 
\begin{equation}
 \mathcal H _{\textup{rat}}
 = \textup{clos}\{f\in L^2(X)\colon  
 f = T^k f, \textup{for some integer $k\geq 1$}   \}, 
 \end{equation}
and $\mathcal H _{\textup{0}}$ is the closed orthogonal space.  
These spaces are invariant under action of $T$.  


Let $\mathbf{1}_{\tilde A} = f_r + f_0$ be the corresponding decomposition of $\mathbf{1}_{\tilde A}$. Now $f_r \in \mathcal H _{\textup{rat}}$, and 
constants are in  $\mathcal H _{\textup{rat}}$, hence  $\lVert f_r\rVert_2^2 \geq  \mu(\tilde A)^2$. 
Moreover, for any integer $n$, we have 
\begin{equation}
    \int_X \mathbf{1}_{\tilde A}  T^n \mathbf{1}_{\tilde A}  \; d\mu 
    =
    \int_X f_r T^n f_r \;d\mu 
    +\int_X f_0 T^n f_0 \;d\mu . 
\end{equation}

 
Thus, given $0< \epsilon < \mu(\tilde A) ^2/5 $,  we can choose $h\in \mathcal H _{\textup{rat}}$,  so that 
$\lVert f_r - h \rVert_2 < \epsilon $. 
And also there is an integer $t$ so that  $h = T^t h$. 
In particular, it follows that $\int_X h \cdot T^{jt} h \; d\mu =\lVert h\rVert_2^2\geq 
\mu(\tilde A)^2- 2\epsilon $ for any $j$. 
Moreover, as $\lVert f_r\rVert_2\leq 1$, we have $\bigl\lvert\int_X f_r\cdot T^mf_r\;d\mu-\int_X h\cdot T^mh\;d\mu\bigr\rvert\leq 3\epsilon$ for every integer $m$.  



We first address the limit obtained by replacing the indicator by $h$.  
The polynomial $p$ is integer valued, so its coefficients are rational; let $L\geq1$ be a common denominator of them, and set $t'=tL$. 
Let $B$ be the set of integers $n$ 
such that  $k_n\equiv 0 \mod t'$.  
This set has positive density by Lemma \ref{lem:positiveDensity}. 
Since the polynomial has no constant term, writing $p(x)=\sum_{i\geq1}c_ix^i$ with $Lc_i\in\mathbb{Z}$, for $k=t'v$ we have 
$c_ik^i=(Lc_i)\,t^i L^{i-1}v^i\in t\mathbb{Z}$. 
Hence $p(k_n)\equiv 0\mod t$ for all $n\in B$, and   
 $h = T ^{p(k_n)}h$.  
(Note that $t\mid k_n$ alone would not suffice: for $p(x)=x(x+1)/2$ and $t=2$, we have $p(2)=3$.) 
That is, only forming the average over the set $B$, we have 
\begin{align} \label{e:relPrime}
\liminf _{N} \Expectation _{\substack{ n\in B\cap  [N] \\  }} \int_X 
f_r \cdot T ^{p(k_n)} f_r \;d\mu  
\geq   \mu(\tilde A) ^2 - 5\epsilon . 
\end{align}
To conclude this case, it remains to argue that the function $f_0$ satisfies 
\begin{align} 
\lim _{N} \Expectation _{\substack{ n\in B\cap  [N] \\  }} \int_X 
f_0 \cdot T ^{p(k_n)} f_0 \;d\mu  
=0. 
\end{align}
And, since $B$ has positive density, it suffices to show that 
\begin{align} 
\lim _{N} \Expectation _{\substack{ n\in [N] \\  }} 
\mathbf{1}_{\{ k_n\equiv 0 \mod t'\} }
\int_X 
f_0 \cdot T ^{p(k_n)} f_0 \;d\mu  
=0. 
\end{align}
Since $\mathbf{1}_{ k_n\equiv 0 \mod t'}=\Expectation_{0\leq a<t'}e(\frac{a}{t'}k_n)$, 
it suffices to show that for each $0\leq a < t'$, 
\begin{align}  \label{eExpPoy}
\lim _{N} \Expectation _{\substack{ n\in  [N] \\  }} \int_X e\left(\frac{a}{t'}k_n \right) 
f_0 \cdot T ^{p(k_n)} f_0 \;d\mu  
=0.
\end{align}
This is a variant of our  polynomial Theorem \ref{t:polyUniform}.  
We address it via spectral methods.   
We know that 
\begin{align} 
\lim _{N} \Expectation _{\substack{ n\in  [N] \\  }} e\left(\frac{a}{t'} k_n  + \alpha p(k_n) \right) =0
\end{align}
for all irrational $\alpha$, by Theorem \ref{t:polyUniform}, since $p$ is non-constant.  
Letting $ \nu   $ be the Herglotz measure for the sequence $\int_X f_0 T^n f_0\;d\mu$, the expectation in \eqref{eExpPoy} is then 
\begin{equation}
   \int_{\mathbf{T}}\Expectation _{\substack{ n\in  [N] \\  }}  e\left(\frac{a}{t'} k_n  + \alpha p(k_n) \right)  \;d \nu (\alpha)  
\end{equation}
As $T$ is invertible, its Koopman operator is unitary, and $\nu$ is the spectral measure of $f_0$. 
An atom of $\nu$ at a rational $a/k$ would produce a non-zero $g$ with $Tg=e(a/k)g$, lying in the closed span of $\{T^nf_0\colon n\in \mathbb{Z}\}\subset \mathcal H_{\textup{0}}$. 
But then $g=T^kg$, so $g\in \mathcal H_{\textup{rat}}$, a contradiction. 
Thus $\nu$ has no rational point masses, and as $\mathbb{Q}/\mathbb{Z}$ is countable, $\nu(\mathbb{Q}/\mathbb{Z})=0$. 
Hence, by dominated convergence, the limit above is zero.  

Combining these, the averages of $\int_X \mathbf{1}_{\tilde A}\cdot T^{p(k_n)}\mathbf{1}_{\tilde A}\;d\mu$ over $n\in B\cap[N]$ have lower limit at least $\mu(\tilde A)^2-5\epsilon>0$. 
Since these integrals are non-negative, and $B$ has positive density $\pi$, the lower limit in \eqref{eENA} is at least $\pi(\mu(\tilde A)^2-5\epsilon)>0$. 
This proves \eqref{e;CantorDensity}, and the final assertion of the Theorem follows since the quantities $d^\ast(A\cap (A-p(k_n)))$ lie in $[0,1]$. 

It remains to see that $H_p$ is intersective. 
Observe first that the argument above establishes \eqref{eENA} for an \emph{arbitrary} invertible measure preserving system $(X,\mathcal B,\mu ,T)$ and any set $\tilde A$ with $\mu (\tilde A)>0$; 
the correspondence principle was used only to produce such a system. 
Observe also that $Z=\{n\colon p(k_n)=0\}$ has at most $\deg p$ elements, as $p$ is non-constant, and that $\lvert p(k_n)\rvert \in H_p$ for every $n\notin Z$. 

Suppose that $I(H_p,N)\not\to 0$. 
The quantity $\Lambda (N)=\max \{\lvert A\rvert \colon A\subset [N], \ (A-A)\cap H_p=\emptyset \}$ is subadditive, since for $A\subset [M+N]$ the traces of $A$ on $[M]$ and on $(M,M+N]$ are each admissible. 
Hence $I(H_p,N)=\Lambda (N)/N$ converges, necessarily to some $\lambda >0$. 
For each $N$, fix $A_N\subset [N]$ with $\lvert A_N\rvert \geq \lambda N$ and $(A_N-A_N)\cap H_p=\emptyset$, and let $\omega_N\in \Omega \coloneqq \{0,1\}^{\mathbb{Z}}$ be the indicator function of $A_N$. 
Let $\sigma $ denote the shift on the compact metric space $\Omega $, let $C=\{\omega \in \Omega \colon \omega (0)=1\}$, and set $\varrho_N=\frac 1N\sum_{n=1}^{N}\delta_{\sigma ^n\omega_N}$, a probability measure on $\Omega $. 
Then $\varrho_N(C)=\lvert A_N\rvert /N\geq \lambda $, while $\varrho_N(C\cap \sigma ^{-h}C)=0$ for every $h\in H_p$, as $A_N$ has no two elements differing by $h$. 
Let $\varrho $ be a weak$^\ast$ limit point of the $\varrho_N$; it is $\sigma $-invariant, and as $C$ and $C\cap \sigma ^{-h}C$ are clopen, 
$\varrho (C)\geq \lambda $ and $\varrho (C\cap \sigma ^{-h}C)=0$ for all $h\in H_p$. 
Since $\sigma $ is invertible, $\varrho (C\cap \sigma ^{-m}C)=0$ for every integer $m$ with $\lvert m\rvert \in H_p$, in particular for $m=p(k_n)$ with $n\notin Z$. 
As $Z$ is finite, the left side of \eqref{eENA}, computed in the system $(\Omega ,\varrho ,\sigma )$ with $\tilde A=C$, is therefore $0$. This contradicts \eqref{eENA}. 
That completes our proof. 
\end{proof}




